% ****************************************************************************************************
% Capítulo 8: O Tempo e o Preço
% ****************************************************************************************************
\chapter{O Tempo e o Preço}
\label{ch:tempo-preco}

% ****************************************************************************************************
\section{A Fila Invisível}
\label{sec:fila-invisivel}
% ****************************************************************************************************

Há uma cena que se repete diariamente em consultórios psiquiátricos ao redor do mundo. Uma pessoa --- que talvez tenha levado meses para superar o estigma, semanas para conseguir uma consulta, anos para admitir que precisava de ajuda --- finalmente senta-se diante do psiquiatra. Conta sua história. Descreve seu sofrimento. E então recebe uma receita acompanhada de uma frase que se tornou tão rotineira que ninguém mais a questiona: \enquote{Este medicamento leva algumas semanas para fazer efeito. Tenha paciência.}

Paciência. A pessoa que chegou em crise, que não consegue trabalhar, que não consegue dormir, que talvez não consiga ver sentido em continuar vivendo --- essa pessoa é instruída a ter paciência. A aguardar. A entrar em uma fila invisível onde seu sofrimento será processado no tempo da molécula, não no tempo da urgência.

Este capítulo examina duas dimensões frequentemente negligenciadas da psicofarmacologia contemporânea: a \textit{temporalidade} do tratamento e o \textit{preço} que se paga por ele. Argumenta-se que ambas as dimensões são tratadas como dados naturais --- propriedades intrínsecas das moléculas --- quando na verdade refletem escolhas, valores e hierarquias implícitas que raramente são discutidas com quem vai vivê-las.

% ****************************************************************************************************
\section{O Tempo do Tratamento e o Tempo do Sofrimento}
\label{sec:tempo-tratamento}
% ****************************************************************************************************

\subsection{A Latência Como Dogma}

A afirmação de que antidepressivos levam duas a quatro semanas para \enquote{começar a fazer efeito} adquiriu, na prática clínica, o estatuto de verdade inquestionável. Como observam Machado-Vieira et al. (2023), \enquote{antidepressivos convencionais que atuam em receptores monoaminérgicos requerem várias semanas para serem eficazes. Este atraso representa um problema significativo nos tratamentos atualmente disponíveis para depressão grave.}

Um problema significativo. A linguagem é curiosamente eufemística. O que significa, concretamente, esse \enquote{problema}? Significa que uma pessoa em episódio depressivo maior --- condição associada a risco substancial de suicídio --- receberá um tratamento que não oferece alívio imediato. Significa que alguém cuja capacidade funcional está gravemente comprometida permanecerá comprometido por semanas adicionais.

E no entanto, essa latência é tratada como propriedade farmacológica inevitável, como se fosse lei da natureza. Mas será?

\subsection{A Prova de Que Poderia Ser Diferente}

Em 2000, Berman e colaboradores publicaram um estudo que deveria ter revolucionado a psiquiatria: uma dose subanestésica de ketamina produziu melhora significativa dos sintomas depressivos em \textit{horas}, não semanas. Em 2019, a FDA aprovou a esketamina intranasal para depressão resistente --- reconhecimento regulatório de que antidepressivos de ação rápida não eram apenas possíveis, mas clinicamente viáveis.

A existência da ketamina como antidepressivo de ação rápida levanta uma pergunta que a psiquiatria prefere não formular diretamente: \textit{se resposta rápida é possível, por que tratamentos de ação lenta permanecem como primeira linha?}

As respostas habitualmente oferecidas --- custo, logística de administração, potencial de abuso --- são válidas, mas incompletas. Não explicam por que a urgência do paciente não é tratada como urgência pelo sistema.

\subsection{O Paciente Como Request em Fila}

O paciente psiquiátrico é tratado como \foreign{request} em uma fila de processamento. Ele entra no sistema, recebe um número (um diagnóstico), é encaminhado para processamento (tratamento), e aguarda sua vez de ser atendido. O tempo de espera é propriedade do sistema, não urgência do solicitante.

O que significa, fenomenologicamente, \textit{esperar} quando se está em crise? Minkowski (1933/1970), em seus estudos sobre tempo vivido, demonstrou que estados depressivos alteram profundamente a experiência temporal. O futuro se fecha, o presente se arrasta, cada minuto pesa como hora. Dizer a alguém nesse estado que \enquote{espere algumas semanas} é impor uma temporalidade objetiva sobre uma temporalidade subjetiva radicalmente distinta.

% ****************************************************************************************************
\section{O Preço Não Anunciado}
\label{sec:preco-nao-anunciado}
% ****************************************************************************************************

\subsection{Emotional Blunting: Efeito Colateral ou Mecanismo de Ação?}

Se a temporalidade do tratamento é raramente questionada, o \textit{preço} do tratamento é ainda mais silenciado. Não me refiro ao preço financeiro, mas ao preço existencial: o que se perde quando se ganha alívio sintomático.

O fenômeno do \foreign{emotional blunting} --- embotamento emocional --- ilustra essa dimensão com clareza perturbadora. Pacientes em uso de ISRSs frequentemente relatam uma espécie de anestesia afetiva: redução não apenas da tristeza patológica, mas da capacidade de experimentar todo o espectro emocional. Alegria, excitação, raiva, ternura --- tudo se atenua. A pessoa deixa de ser inundada pela dor, mas também deixa de ser tocada pela vida.

A pergunta que emerge é incômoda: o emotional blunting é \textit{efeito colateral} dos ISRSs ou \textit{mecanismo de ação}? Se a redução da ansiedade patológica vem acompanhada de redução da capacidade de sentir em geral, estamos \textit{tratando} a ansiedade ou \textit{suprimindo} a afetividade?

\subsection{A Disfunção Sexual Como Regra}

A disfunção sexual induzida por antidepressivos é talvez o exemplo mais documentado de preço não anunciado. Estudos estimam que 50\% a 70\% dos pacientes em ISRSs experimentam alguma forma de comprometimento sexual.

A magnitude desses números deveria provocar reflexão. Não estamos falando de efeito raro. Estamos falando de \textit{maioria}. A disfunção sexual é, para ISRSs, mais regra que exceção.

Mais perturbador ainda é o fenômeno da \foreign{Post-SSRI Sexual Dysfunction} (PSSD) --- disfunção sexual que persiste após a descontinuação do medicamento, potencialmente de forma permanente. A Agência Europeia de Medicamentos reconheceu o PSSD em 2019, exigindo atualização das bulas.

\subsection{O Que Não É Dito}

Retornemos à cena do consultório. O psiquiatra prescreve um ISRS. O que é dito ao paciente?

Tipicamente: que o medicamento ajudará com a ansiedade ou depressão; que leva algumas semanas para fazer efeito; que pode haver efeitos colaterais no início, geralmente transitórios.

O que frequentemente \textit{não} é dito: que há probabilidade de 50--70\% de alguma disfunção sexual; que essa disfunção pode persistir após parar o medicamento; que o paciente pode experimentar embotamento emocional; que existem alternativas com perfil diferente de efeitos colaterais.

Por que esse silêncio? Parte da resposta é pragmática: médicos temem que informação detalhada comprometa a adesão. Mas essa justificativa é paternalista --- assume que o paciente não é capaz de fazer escolhas informadas sobre seu próprio corpo.

% ****************************************************************************************************
\section{Quem Decide O Que Vale O Quê?}
\label{sec:quem-decide}
% ****************************************************************************************************

\subsection{A Hierarquia Implícita de Valores}

Toda prescrição psiquiátrica implica uma hierarquia de valores. Quando prescrevo um ISRS sabendo que provavelmente causará disfunção sexual, estou afirmando --- implicitamente, sem consultar o paciente --- que \textit{não ter ansiedade} vale mais que \textit{ter vida sexual satisfatória}.

Essas hierarquias não são neutras. Refletem valores culturais, vieses profissionais, interesses institucionais. Refletem também, com frequência perturbadora, a perspectiva de quem prescreve mais que a perspectiva de quem vai viver com as consequências.

Para um psiquiatra, o sucesso do tratamento é medido por redução de sintomas em escalas padronizadas. A disfunção sexual não aparece na escala de Hamilton para depressão. O embotamento emocional não é critério de remissão. São, literalmente, invisíveis para a métrica que define \enquote{melhora}.

Mas para o paciente, esses \enquote{efeitos colaterais} podem ser centrais à experiência de estar vivo.

\subsection{A Ausência de Decisão Compartilhada}

O conceito de \foreign{decisão compartilhada} tornou-se mantra na medicina contemporânea. Na psiquiatria, essa prática permanece mais ideal que realidade. A estrutura da consulta de 15 minutos não permite discussão aprofundada de trade-offs. O paciente recebe a receita; não participa da escolha que a gerou.

\subsection{O Paciente Como Objeto}

Há uma inversão epistemológica em jogo. O paciente --- que veio ao consultório como sujeito de seu sofrimento, buscando tornar-se sujeito de sua cura --- é transformado em \textit{objeto} da prescrição. Recebe tratamento como se recebe procedimento: passivamente, sem participação na definição de meios ou fins.

O silêncio sobre trade-offs não é apenas omissão; é violação. Violação do direito à informação, do direito à autonomia, do direito de ser tratado como pessoa capaz de fazer escolhas sobre a própria vida.

% ****************************************************************************************************
\section{Uma Farmacologia Centrada no Sujeito}
\label{sec:farmacologia-sujeito}
% ****************************************************************************************************

\subsection{O Tempo do Paciente}

Uma farmacologia centrada no sujeito começaria por reconhecer que o tempo do tratamento deve dialogar com o tempo do sofrimento. Isso implicaria: avaliação explícita do risco durante o período de latência; estratégias de suporte intensivo nas primeiras semanas; consideração de tratamentos de ação rápida para casos de maior urgência.

\subsection{Os Valores do Paciente}

Uma farmacologia centrada no sujeito perguntaria, antes de prescrever, quais valores o paciente prioriza. Para alguns, alívio da ansiedade a qualquer custo é prioritário. Para outros, preservar a vida sexual é inegociável. Não há resposta universal correta --- há respostas singulares que dependem de quem vai viver com as consequências.

Isso exige tempo. Exige conversa. Exige o que a psiquiatria contemporânea frequentemente não oferece: encontro intersubjetivo em que médico e paciente, juntos, deliberam sobre o que fazer.

\subsection{O Preço Explicitado}

Uma farmacologia centrada no sujeito explicita os preços antes de cobrar. Informa sobre disfunção sexual antes de prescrever ISRS. Discute embotamento emocional como possibilidade real. Apresenta alternativas com perfis diferentes de benefícios e custos. Permite que o paciente escolha qual preço está disposto a pagar.

% ****************************************************************************************************
\section{A Urgência e o Custo}
\label{sec:urgencia-custo}
% ****************************************************************************************************

Iniciamos este capítulo com uma cena de consultório --- o paciente que recebe prescrição e instrução para ter paciência. Percorremos a temporalidade do tratamento e os trade-offs que o acompanham. A conclusão que emerge é simples: a psicofarmacologia contemporânea opera com um tempo e um sistema de valores que frequentemente não são os do paciente.

O tempo do tratamento é o tempo da molécula --- semanas para receptores se regularem. O tempo do sofrimento é outro --- urgente, insuportável, medido em minutos que não passam. Tratar o primeiro como dado natural e o segundo como variável a ser \enquote{manejada} é inverter prioridades.

O preço do tratamento é definido pelo prescritor --- efeitos colaterais são \enquote{toleráveis} quando quem prescreve não vai tolerá-los. Quem decide que não ter ansiedade vale mais que ter vida sexual? Se o paciente não participa dessas decisões, não está sendo tratado como sujeito; está sendo tratado como objeto de intervenção farmacológica.

% ****************************************************************************************************
% Referências do Capítulo
% ****************************************************************************************************
\vspace{2em}
\begin{small}
\noindent\textsc{Referências}

\vspace{1em}
\noindent Berman, R. M., et al. (2000). Antidepressant effects of ketamine in depressed patients. \textit{Biological Psychiatry}, 47(4), 351--354.

\noindent European Medicines Agency. (2019). PRAC recommendations on signals: Adopted at the 10-13 June 2019 PRAC meeting.

\noindent Machado-Vieira, R., et al. (2023). Ketamine and the next generation of antidepressants. \textit{Pharmacology \& Therapeutics}.

\noindent Minkowski, E. (1970). \textit{Lived time}. (Orig. 1933). Northwestern University Press.

\noindent Montejo, A. L., et al. (2001). Incidence of sexual dysfunction associated with antidepressant agents. \textit{Journal of Clinical Psychiatry}, 62(Suppl 3), 10--21.

\noindent Price, J., et al. (2023). Emotional blunting in depression: The role of serotonin. \textit{Neuropsychopharmacology}.
\end{small}
